\section{Main results and how to read them}\label{sec:overview}

The answer is positive for a broad, precisely specified core of category
theory. Its structural arguments use very little arithmetic: the real work
is to give their data a legitimate arithmetic interface. One can permit
arithmetic, potentially undecidable equality of arrows without allowing
arbitrary function spaces or induction on all arithmetic formulas.

This report uses the user's overview \cite{wiki} to delimit the basic
subject, and standard author-hosted references for terminology
\cite{riehl,awodey}. Its mathematical arguments are written out in full.
The classical theorems themselves are not claimed to be new. The purpose
is a coherent iPRA realization with every computational and logical
boundary visible, building on the pinned ProveIt development.

\subsection{The principal theorem families}
\begin{longtable}{@{}>{\raggedright\arraybackslash}p{0.27\linewidth}>{\raggedright\arraybackslash}p{0.60\linewidth}p{0.06\linewidth}@{}}
\toprule
Result & Exact form established here & Sec.\\
\midrule
\endhead
Arithmetic foundations & Fixed formulas for objects, arrows and equality;
PR code operations; specialization of finite structural proofs. &
\ref{sec:foundation}\\
Functorial calculus & Functors, natural transformations, whiskering,
horizontal and vertical composition, interchange, duality. &
\ref{sec:functors-yoneda-foundation}\\
Yoneda & Explicit mutually inverse constructions
$\alpha\mapsto\alpha_A(\id_A)$ and $x\mapsto(f\mapsto P(f)x)$,
with all typing and naturality laws. & \ref{sec:yoneda}\\
Representability & Full faithfulness, canonical uniqueness of representing
objects, and representation by a chosen universal element. &
\ref{sec:yoneda}\\
Density / co-Yoneda & Canonical generator normal form and a universal
cocone, with arithmetic equality defined by equality of values. &
\ref{sec:density}\\
Universal constructions & Unique isomorphisms; pullbacks; all fixed finite
limits from selected terminal objects, products and equalizers, plus duals. &
\ref{sec:limits}\\
Finite-set existence & Uniform PR algorithms for finite limits, colimits
and exponentials, including $[0]$ and $0^0=1$. & \ref{sec:finite}\\
Adjunctions & Hom correspondences iff unit/counit triangles; uniqueness;
preservation of supplied limits and colimits. & \ref{sec:adjunctions}\\
Equivalence & A quasi-inverse constructed from faithful, split full and
split essentially surjective data. & \ref{sec:adjunctions}\\
Monads & Monad of an adjunction; Kleisli and Eilenberg--Moore categories
and adjunctions; creation of supplied algebra limits. & \ref{sec:monads}\\
Idempotents & Canonical splitting in the Karoubi envelope and extension
along displayed splittings. & \ref{sec:limits}\\
Kan extensions & Pointwise formula with supplied uniform comma
(co)limits; unconditional existence for finite tables in finite sets. &
\ref{sec:kan}\\
Diagonal theorem & Lawvere fixed points and the obstruction to an
unrestricted PR evaluation family. & \ref{sec:lawvere}\\
\bottomrule
\end{longtable}

\subsection{Three levels of assertion}
A result about a fixed presentation is an ordinary arithmetic theorem.
A result quantified over a numerical family of finite tables is also an
ordinary arithmetic theorem, now with table codes as variables.
A result ranging over arbitrary presentation syntax is an external
schema producing such theorems. These levels are distinguished throughout.

The structural theorems are relative to the explicitly stated category
laws and selected universal operations. When those laws have iPRA proofs,
the theorem's proof specializes to an iPRA proof. This is mathematically
substantial: the Yoneda and adjunction calculations apply also to infinite
presented categories with undecidable arrow equality. They are not
restricted to testing finite tables.

Conversely, an existence claim that a category has all limits
requires more than a formal calculation. The finite-set section supplies
actual constructors and a quantifier-free induction analysis. For
arbitrary presentations, the report states the exact universal operations
needed and does not invent them.

\subsection{What the deliverable's verification labels mean}
\begin{description}[style=nextline]
\item[Human proof]
The displayed assumptions and conclusions are connected by complete
mathematical arguments in this report. Numerical algorithms include
their invariants and primitive-recursive bounds; structural proofs use
only finite intuitionistic inferences.
\item[Reused object-language source]
The pinned ProveIt modules actually construct judgments of the form
$\Gamma\vdash[\mathsf{PRAAx}]\varphi$. They were read and their import
closure included. The repository reports its derivation library as
kernel-checked; it was not rebuilt in this environment.
\item[New Lean integration draft]
The three new modules target that same derivation syntax and contain
no proof placeholders, but were not compiled because the Lean executable
was unavailable. Their elaboration and kernel checking remain to be checked.
\item[Exact finite-model check]
A reproducible Python run checks the stated finite cases with exact
integer arithmetic. It supports implementation validation but is not
an iPRA formalization.
\end{description}

Useful refinements relative to the inspected upstream sources include
an interface for Yoneda at one fixed parameter and representing object;
a direct density presentation that avoids arbitrary quotient existence;
a complete empty-set treatment and explicit uniform finite-diagram
tester; and one coherent proof and formalization plan for the monad,
algebra, Kan-extension and idempotent constructions.
No claim of a completed Lean mechanization of all these refinements is made.

